\begin{afterword}
\summaryhead

The post describes the bystander effect: larger groups are less likely to act in an emergency, even taken together. In Latané and Darley's smoke experiment, 75\% of subjects alone reported smoke entering the room, but only 38\% of groups of three did, and only 10\% of subjects sitting with two confederates who ignored it. The standard explanations are diffusion of responsibility and pluralistic ignorance, which the post illustrates with a passage from Cialdini. Cialdini advises a person in need to single out one bystander.

The author considers an evolutionary arms race to be the last to help and rejects it, since the experiments show a failure to help rather than a delay. The author suggests instead that the effect arises because strangers are a modern problem, since helping rises when people know the victim or each other, and that fear of seeming to bid for status may play a part. Finally the post argues that bystanders are not merely avoiding blame, because they ignored a threat to themselves and because teaching people about the effect reduces it.

\responsehead

The post reports its experiments accurately, and it marks its speculations as speculations. Several of its guesses were already made or supported in the article it cites.

The figures for the smoke study match Latané and Darley's article, and Cialdini's advice is reported correctly. Two cautions belong beside them. The samples are small: 24 subjects alone, ten with confederates, eight groups of three. And the claim that groups fail ``collectively'' is shown for the laboratory studies in that article. A 2020 study of real public conflicts filmed by CCTV found that someone intervened in nine of ten cases, and that more bystanders went with a greater chance that someone would. That study came eleven years after the post; it bears on how far the opening sentence can be generalized, not on what the author could have known.

The cited article opens with the murder of Kitty Genovese and the thirty-eight neighbors said to have watched in silence. The post does not repeat the story, and its argument does not need it. The reader should still know that the story was later shown to be largely false. A 2007 study of the archives found no evidence for 38 witnesses, for their seeing the murder, or for their inaction, and in 2016 the New York Times said its 1964 report had grossly exaggerated the number of witnesses and what they had perceived.

The speculative middle of the post is careful. The author raises the guess that people evolved to wait for someone else to act first, a guess the author had made in ``Evolving to Extinction'' in 2007, and rejects it on a premise the seizure study supports: subjects who had not responded within three minutes never did. Both findings offered for the replacement guess, that knowing the victim or each other raises helping, are in the cited article. The author cannot recall shyness being discussed, and blames ``poor memory''; the article discusses fear of embarrassment before an audience, and uses it to explain why friends inhibit each other less.

The claim that teaching people about the effect reduces it comes without a source. A 1978 study supports it. The post's conclusion, that bystanders are not merely avoiding blame, agrees with Latané and Darley, who found that non-responders were in conflict, not indifferent.

In short: an accurate report of Latané and Darley's experiments, with hedged guesses that the cited article had in several cases already made, and a claim of collective failure based on small laboratory studies.
\end{afterword}
